\addcontentsline{toc}{chapter}{ABSTRACT}
\begin{center}
  {\bfseries\fontsize{16pt}{19pt}\selectfont ABSTRACT}
\end{center}
\vspace{0.8cm}

\noindent
Quantum key distribution (QKD) networks require routing decisions that account for
finite link-local key resources and changing conditions. This thesis presents
\QFlow, a deterministic fixed-point FPGA accelerator for evaluating externally
generated route candidates in trusted-relay QKD networks. The upstream KMS/SDN
control plane generates candidates, evaluates hard constraints, and supplies each
candidate's validity bit and bounded descriptors; authorization and route
installation also remain outside the FPGA.

Four features are quantized into 256 states with four profile actions. Tabular
Q-learning is performed offline, and the mapping is deployed as a 256-entry two-bit
policy ROM with minimum-dwell control. Training uses four coefficients to construct
profile-dependent path cost and three Tchebycheff ratios. During physical replay,
path cost is precomputed and common across profiles, so H2--H4 deploy only the ratio
projection. Physical exactness therefore validates that projection, not on-chip
execution of the complete training action. Five independent configurations comprise
a feasible minimum-hop baseline (H0), fixed key-aware baseline (H1), fixed-profile
\QFlow\ (H2), rule-adaptive selection (H3), and learned profile selection (H4).

The experiment uses minimum key occupancy, bounded route quality, offered load, and
utilization imbalance. Route quality is an abstract input, not quantum coherence of
post-processed classical keys. The network model separately represents key-pool
generation, consumption, status, utilization, freshness policy, and health.

Verification combined a fixed-point reference, exhaustive tests, cycle-accurate RTL
simulation, independent Xilinx ISE implementations, and physical replay on a Nexys 3
Spartan-6 FPGA. Across five builds, 420 records and 3,780 checked fields produced zero
mismatches; every configuration exceeded 100~MHz. H4 achieved 102.020~MHz and mean
accepted-request-to-done kernel latency of approximately 2.048~$\mu$s. Relative to
H3, H4 added 0.986\% mean latency, reduced maximum frequency by 0.612\%, and produced
24 rather than 44 profile transitions, with fewer LUTs but more registers and
occupied slices. Reproducible paired-row confidence intervals crossed zero, with
sign-test $p$-values 0.2188 and 0.4531; route-quality superiority was not established.

Supporting synthesis and OpenROAD results concern isolated kernels, not an
integrated ASIC. The contribution is a bit-exact, microsecond-scale adaptive
candidate evaluator with explicit semantic and measurement limits.

\vspace{0.5cm}
\noindent\textbf{Keywords:} quantum key distribution network, FPGA accelerator,
route-candidate evaluation, fixed-point arithmetic, reinforcement learning,
policy ROM, trusted relay.
